\section{Definition 1}\label{s:def}

\lem{2}{A life $L$ is \emph{meaningful to degree $d$} if $d$ is the extent to which $L$ consists in active engagement with projects of objective worth \citep{Wolf2010}.}

\paragraph{The sentence as a whole.} Definition~1 defines a two-place predicate, \enquote{meaningful to degree $d$}, true of a life and a degree. The definiendum is italicized, as is customary, and the definiens follows \enquote{if}. The definiens has three components: a bearer (a life), a measure (the extent to which the life consists in something) and a condition (active engagement with projects of objective worth). Each is the subject of a phrase lemma below.

Three features of the sentence as a whole should be noted before the phrases are examined. First, the definition mentions no subject. It does not say that a life is meaningful to the extent that \emph{its subject} is engaged, as EXI-17-a (Principle~P1) and EXI-11-a (Definition~2.1) do; it says that a life is meaningful to the extent that it \emph{consists in} engagement. The subject is implied, since engagement is someone's engagement and a life is someone's life, but the grammatical subject of the condition is the life. The choice is deliberate and has a consequence examined at \Lm{2.4}: it makes meaning a matter of what the life is made of, not of what the subject brings about or of how she regards her life.

Second, the definition is a real definition, not a stipulation. It is cited to \citet{Wolf2010}, whose lectures argue that the concept we use when we call lives meaningful or meaningless is the concept of a life engaged with worthwhile things. A stipulative definition would need no citation and could not be wrong. Definition~1 can be wrong: if there were lives that reflection judged meaningful and that contained no active engagement with anything of objective worth, or lives so engaged that reflection judged meaningless, the definition would fail. The comment on \Lm{2.7} considers the cases that have been offered.

Third, the definition is general: it applies to every life, actual or possible. This matters for Theorem~4, which concerns the lives available to an agent and not only the life she actually leads. A definition that applied only to actual lives could not support a claim about availability, which is a claim about possible lives (\Lm{5.3}).

\paragraph{Parallels in the collection.} The definition is stated three times in the discipline, in slightly different words. EXI-17-a states it as a principle: \enquote{A life is meaningful to the extent that its subject is actively engaged with what is worthwhile} (its Principle~P1). EXI-11-a states it as a definition with the same words (its Definition~2.1). EXI-01-c restates Definition~1 and takes it over without change. The differences between EXI-01-a and the later statements are three: the later statements mention the subject, they replace \enquote{projects} by \enquote{what is worthwhile}, and they drop the explicit degree variable while keeping \enquote{to the extent that}. The second difference is the most important, and the comment on \Lm{2.6} argues that it is a clarification of EXI-01-a rather than a departure from it. The concordance (Appendix~\ref{app:conc}, Table~A.2) shows how the vocabulary of the definition is distributed across the discipline: the families of \emph{engage-} and \emph{worth-} occur \cnFamengag{} and \cnFamworth{} times in the other manuscripts, against \cnFamproject{} occurrences of \emph{project}, which is the measure of how far the later manuscripts have moved from the word \enquote{project} to the more general vocabulary of objects of engagement.

\lem{2.1}{A life $L$}

The bearer of meaning is a life. Three alternatives are excluded by this choice.

The first alternative is the person. One speaks of meaningful lives, not of meaningful persons, and the usage reflects a real distinction: a person's character or worth is assessed at a time, while a life is assessed as a temporally extended course. A person whose life has been empty may become someone whose life is meaningful; we do not say that she has become a meaningful person. Definition~1 follows usage. The second alternative is the action or the episode. Actions and episodes can be meaningful in a sense: a gesture can be meaningful, and a conversation can be the most meaningful hour of a week. Definition~1 does not deny this, but it does not define it. It defines the meaning of a life, and the meaning of the parts of a life enters only through the measure (\Lm{2.3}). The third alternative is humanity or life as such, which the comment on \Lm{1} has set aside.

\paragraph{Whole and part.} Metz distinguishes the meaning of a life as a whole from the meaning in a life, the meaning of its parts, and asks whether the first is the sum of the second or depends also on the relations among the parts \citep{Metz2013}. \citet{Kauppinen2012} argues that the meaning of a life depends on its narrative structure, on whether its later parts build on and redeem its earlier ones. \citet{MacIntyre1981} had earlier argued that a human life has the unity of a narrative and that its goods are intelligible only within that unity. Definition~1 takes a position on this question only through the phrase \enquote{the extent to which $L$ consists in}: the degree of meaning of the whole is the extent to which the whole consists in a certain kind of activity. On its most natural reading, this makes whole-life meaning depend on the meaning of the parts, without saying how. It is compatible with the view that the parts contribute additively and with the view that their arrangement matters. EXI-19 holds that the shape of a life, the order of its goods and evils, bears on its value as a whole (Principle~P1 of EXI-19-a), and EXI-19-b argues that finitude is a condition of the meaning of a life as a whole (its Proposition~1.2). These claims concern a kind of meaning that Definition~1 neither includes nor excludes; Volume~III examines whether they require a revision of the measure.

\paragraph{Lives as possibilities.} The variable $L$ ranges over lives, and Theorem~4 quantifies over lives available to an agent. A life in this sense is a possible course of an agent's existence, individuated by its subject and its content. The Lean module makes this explicit: \texttt{Life} is a family of types indexed by agents, so that each agent has her own space of possible lives (Section~\ref{s:lean}). Two features of this conception should be noted. A life is someone's life; there are no lives without subjects. And a life is a whole; the life of an agent who dies young is a complete life, not a fragment of the longer life she might have had.

\lem{2.2}{is \emph{meaningful to degree $d$}}

\paragraph{Gradualism.} The definiendum makes meaning a matter of degree. This is the common view in the literature: \citet{Wolf1997} speaks of lives as meaningful insofar as they are engaged with worthwhile projects, and \citet{Metz2013} treats meaning as a gradable property throughout. The alternative, that meaning is all or nothing, has few defenders, and Definition~1 does not argue against it. It is nevertheless worth stating why gradualism is right. Reflection compares lives in respect of meaning: the life of someone who devoted thirty years to the care of her disabled brother and the life of someone who gave him an afternoon a month are both meaningful in respect of that care, and the first more so. A binary concept could not register the comparison.

\paragraph{\enquote{Meaningful} simpliciter.} The definition defines \enquote{meaningful to degree $d$} and leaves \enquote{meaningful} without qualification undefined. The text uses \enquote{positive meaning} in Theorem~4 and \enquote{meaningful to a positive degree} in its proof, and Remark~5 uses \enquote{meaningful} without qualification (\Lm{10}). The natural reading is that a life is meaningful simpliciter if and only if it is meaningful to some positive degree. Ordinary usage is more demanding: to call a life \enquote{meaningful} is usually to say that it has a good deal of meaning, as to call a person \enquote{tall} is to say that she is taller than most. On the ordinary usage, Theorem~4 would claim much less than \enquote{a meaningful life is available}, and the family conclusion that meaning is available to everyone would be weaker than it sounds. The text avoids the ambiguity by stating Theorem~4 in terms of positive meaning, and the reader should hold it to that. EXI-01-d makes the point explicit in its reply to the objection that availability so secured is trivial: availability is not possession, and possession comes in degrees.

\paragraph{The range of $d$.} The text does not say what the degrees are. Three possibilities should be distinguished. On the cardinal reading, $d$ is a real number, and differences between degrees are meaningful: one life can be twice as meaningful as another. On the ordinal reading, degrees are ranks in a complete ordering of lives. On the partial reading, degrees are elements of a partial order, so that some pairs of lives are incomparable in meaning. The text is neutral. It uses only one feature of the degrees: that there is a distinguished degree, zero or none, and that degrees above it are positive. Every one of the three readings supplies this.

We think the partial reading is the most plausible, and we record the reason, though nothing in Volume~I depends on it. The meaning of a life spent in scientific research and the meaning of a life spent in the care of others are both substantial, and reflection does not rank them; it does not judge them equal either, since a small improvement in one would not make it more meaningful than the other. This is the mark of incomparability in the sense discussed by \citet{Raz1986}. A cardinal or complete ordinal measure would force a ranking that reflection does not make.

\begin{proposition}[No negative degrees]\label{p:noneg}
On Definition~1, no life is meaningful to a negative degree.
\end{proposition}

\begin{proof}
The degree of a life is the extent to which it consists in active engagement with projects of objective worth. An extent is a measure of how much of something there is, and is not negative. A life that consists in no such engagement is meaningful to the degree of zero, whatever else it consists in.
\end{proof}

Proposition~\ref{p:noneg} has a consequence that some will find objectionable. A life devoted to cruelty, or to the destruction of what is valuable, is on Definition~1 no less meaningful than a life devoted to nothing: both have zero meaning in respect of those activities. One might think that the destroyer's life has negative meaning, that it is not merely empty but anti-meaningful. Definition~1 does not allow this, and we think it right not to. The destroyer's life is bad, and morally worse than the idle life, but the respect in which it is worse is a moral respect. Wolf's argument that meaning is a value distinct from morality \citep{Wolf1997} implies that a life can be bad without being anti-meaningful, and Definition~1 keeps the distinction clean. A theory that allowed negative meaning would have to say what it is about destruction, as opposed to its moral wrongness, that subtracts meaning, and no account of this that does not collapse into moral assessment has been given.

\lem{2.3}{if $d$ is the extent to which}

\paragraph{\enquote{if}.} As is usual in definitions, \enquote{if} is to be read as \enquote{if and only if}. The proof of Theorem~4 uses the definition only in the direction from right to left (from engagement to meaning), and the Lean module formalizes only that direction (Section~\ref{s:lean}). But the definition is a definition, and Remark~5 uses the other direction when it says that a particular life may fail to be meaningful because of what is done with it (\Lm{10}): a life with no engagement with projects of objective worth has no meaning, and that is the left-to-right direction.

\paragraph{\enquote{the extent}.} The definite article implies that each life has exactly one degree of meaning. On the partial reading of the degrees, this is unproblematic: a life has one place in a partial order. What is problematic is the measure. \enquote{The extent to which $L$ consists in} engagement admits two readings that diverge in important cases.

On the \emph{proportional} reading, the extent is the share of the life that consists in engagement: a life of which half is spent in engagement with worthwhile projects has a higher degree than a life of which a quarter is. On the \emph{absolute} reading, the extent is the amount of engagement the life contains: a life with thirty years of engagement has a higher degree than a life with twenty, whatever the length of each.

The readings diverge in cases of the following kind. Two people spend thirty years in devoted engagement with worthwhile work. One dies at sixty; the other lives to ninety, and spends the last thirty years in idleness. On the proportional reading, the first life is more meaningful, since a larger share of it consists in engagement. On the absolute reading, the two lives are equally meaningful, since they contain the same engagement. Reflection is uncertain here, and its uncertainty is instructive. There is something to the thought that the idle years diminish the second life, and something to the thought that they subtract nothing from what it contained.

\begin{proposition}[Measure neutrality]\label{p:neutral}
Theorem~4 holds on every reading of \enquote{the extent to which $L$ consists in} on which a life that contains active engagement with a project of objective worth, sustained over some non-null part of the life, consists in such engagement to a positive extent.
\end{proposition}

\begin{proof}
The proof of Theorem~4 uses Definition~1 only to infer positive meaning from positive extent (\Lm{8}), and the inference from inclusion of engagement to positive extent (\Lm{7.2}). Both the proportional and the absolute readings satisfy the stated condition: a non-null part of a life is a positive share of it, and a positive amount of engagement. Any reading that satisfies the condition supports the same inference.
\end{proof}

Proposition~\ref{p:neutral} explains why the text can leave the measure unspecified: the availability claim does not depend on it. The measure matters for comparisons of degree, and for the claims of EXI-01-d about what affects degree (its Proposition~2.4). Volume~III states a measure and examines its consequences. One feature of the measure can be fixed now: it cannot be purely temporal. Engagement differs in depth as well as in duration, and the worth of its objects differs too, and EXI-11-a states as a principle that engagement with what is of greater worth makes a life more meaningful, other things equal (its Principle~P2). The extent to which a life consists in engagement with projects of objective worth is therefore at least a function of three variables: how much of the life is so engaged, how active the engagement is, and how great the worth of its objects. Definition~1 compresses these into one word. The compression is legitimate in a definition that is meant to say what meaning consists in, and it would be a defect in a theory that claimed to measure it.

\lem{2.4}{$L$ consists in}

\paragraph{Constitution, not causation.} \enquote{Consists in} is the language of constitution. On Definition~1, the meaning of a life is not an effect of engagement, as the satisfactions of a successful career are an effect of the work; it is the engagement, measured. This sets Definition~1 against two families of rival views.

The first family makes meaning depend on consequences. On the good cause account of \citet{Smuts2013}, a life is meaningful to the extent that it causes good outcomes, whether or not its subject is engaged with them and whether or not she knows of them. On the transcendence view of \citet{Nozick1981}, a life gains meaning by making a difference to, and leaving traces in, something of value beyond itself. Both views make the meaning of a life depend on facts outside the life: on what happens as a result of it. Definition~1 makes it depend on what the life consists in.

The difference shows in two kinds of case. In the first, a life produces great goods without engagement: a person invents, idly and as a joke, a device that later saves many lives. On the good cause account, her life is very meaningful. On Definition~1, the invention adds to the meaning of her life only to the extent that her life consisted in engagement with it, which by hypothesis it did not. We think Definition~1 gives the right verdict: her life was fortunate in its effects, and the world owes her gratitude, but her life was not given meaning by an outcome she did not care about. In the second, a life is engaged with a worthwhile project that produces nothing: the unsuccessful mathematician of EXI-17-a, who spends her career on a problem she never solves (its Lemma~3.1). On the good cause account, her life has little meaning in respect of that work. On Definition~1, it has much. EXI-17-a argues that reflection sides with Definition~1, and the argument depends on the constitutive reading of \enquote{consists in}.

\paragraph{Posthumous events.} A consequence of the constitutive reading is that the meaning of a life is fixed by its course and cannot be altered after the life has ended. A painter whose work is discovered after her death did not live a more meaningful life than she would have if the work had been lost; her life consisted in the same engagement in either case. This must be distinguished from the claim of EXI-13 that the dead can be harmed: a setback to a surviving interest is bad for the ante-mortem person (Principle~P1 of EXI-13-a), and the painter's interest in being recognized can be thwarted or satisfied after her death. Harm and meaning come apart here, as they do elsewhere. The painter's life went better for her, perhaps, if her work was recognized; it was not more meaningful.

\paragraph{Meaning without awareness.} A further consequence is that a subject can be mistaken about the meaning of her own life. Whether her life consists in active engagement with projects of objective worth is a fact about her activity and about the worth of its objects, and she can be wrong about the second. Someone who believes that her life is meaningful because she is passionately engaged in a project that has no worth is mistaken; someone who believes that her life is meaningless, because she does not see the worth of the care she gives, is mistaken in the other direction. The possibility of this second kind of mistake is part of what gives the question of meaning its urgency, and it is discussed further under Assumption~2 (Section~\ref{s:defence}).

\paragraph{Meaning and happiness.} EXI-07 holds that happiness is a by-product of other pursuits, so that aiming directly at it is self-defeating; EXI-07-a argues that a large part of happiness consists in or depends on wholehearted engagement in activities and relationships whose objects are valued for their own sake (its Principle~P1). Meaning, on Definition~1, consists in engagement of a certain kind. The two claims fit: engagement with worthwhile objects constitutes meaning and tends to cause happiness. They also explain why meaning and happiness come apart. A life of engagement with worthwhile but painful work, the care of the dying for instance, may be meaningful and unhappy; a life of pleasant engagement with trivial pursuits may be happy and nearly meaningless. Wolf's case for treating meaning as a distinct dimension of the good life rests on such pairs \citep{Wolf1997}.

\lem{2.5}{active engagement}

\paragraph{Wolf's notion.} \enquote{Active engagement} is Wolf's phrase. She describes the engaged subject as gripped, excited, involved, absorbed in what she is doing, as against the subject who goes through the motions, and she contrasts it with the passive enjoyment of comfort or diversion \citep{Wolf2010}. EXI-17-a makes the notion precise: an agent is actively engaged with an object if she acts in ways directed at it, cares how it fares, and is gripped or absorbed by it, so that her activity expresses her concern rather than being merely performed (its Definition~2.1). The definition has three clauses, a practical clause (directed action), a conative clause (care) and an affective clause (absorption), and the commentary on \Lm{2.5} must say how they are related.

\paragraph{Two readings.} On the \emph{affective} reading, engagement is fundamentally a felt state: the experience of being gripped. Action and care matter because they are what the feeling is about. On the \emph{conative} reading, engagement is fundamentally caring about an object in a way that issues in action directed at it. The feeling of absorption is the normal manifestation of such caring, but it is not its essence.

The readings diverge in the case of someone who acts with care but without feeling. A nurse who is clinically depressed continues to work devotedly: she attends to her patients, worries about them, does what she can for them, but she feels nothing in doing so; the work does not grip her, and she takes no satisfaction in it. On the affective reading, she is not actively engaged, and her work adds nothing to the meaning of her life during her illness. On the conative reading, she is engaged, and her depression diminishes her well-being without diminishing the meaning of her life.

We think the conative reading is the better reading of the text and the more plausible view. Three considerations support it. First, the text's adjective is \enquote{active}, not \enquote{felt}: it contrasts engagement with passivity, not with numbness. Second, EXI-17-a's definition ends with the clause \enquote{so that her activity expresses her concern rather than being merely performed}, which makes concern, not feeling, the decisive element. Third, the affective reading would make the meaning of a life depend on a feature of its subject's affective life that is not under her control and that bears no relation to the worth of what she does. The depressed nurse is doing exactly what she would be doing if she were well, for the same reasons and with the same concern; it would be strange if her illness, which takes so much from her, also took the meaning from her work.

The conative reading draws on Frankfurt's account of caring. \citet{Frankfurt1982} argued that what a person cares about is distinct from what she desires or values, that caring consists in a person's identifying herself with what she cares about and making herself vulnerable to its fortunes, and that it gives her life a structure of concern that persists over time. Active engagement on the conative reading is caring in Frankfurt's sense that issues in activity. Frankfurt himself held that the importance of what we care about derives in large part from our caring about it, a view that Definition~1 rejects through the qualifier \enquote{of objective worth} (\Lm{2.7}); but his account of what caring is can be separated from his account of where importance comes from, and Definition~1 takes the first without the second.

\paragraph{Activity and process.} \enquote{Active} also recalls Aristotle's distinction between activity (\emph{energeia}) and process (\emph{kinesis}) (\emph{Metaphysics} $\Theta$.6, 1048b18--35). A process is directed at an end outside itself and is incomplete until the end is reached: building is incomplete until the house is built. An activity is complete at every moment: one who sees has seen, one who contemplates has contemplated. Engagement as Definition~1 understands it has the structure of an activity. The mathematician engaged with her problem is engaged throughout, and her engagement is not incomplete because the problem remains unsolved. This is the structural reason why, as EXI-17-a argues, meaning does not require achievement: achievement belongs to processes, which are completed or not, while the engagement that constitutes meaning is an activity, which is complete in being exercised. The activity can be directed at an end outside itself, as the mathematician's is, without being incomplete until the end is reached.

\paragraph{Appreciation.} \enquote{Active} might seem to exclude receptive engagement: listening to music, contemplating a landscape, reading. It does not. EXI-17-a argues that the listener who spends much of her life studying music she loves, with growing understanding, is actively engaged, since attentive listening and study are activities directed at their object (its Lemma~3.2). Activity in the relevant sense is not productive activity but directed attention, and receptive engagement involves it as fully as production does. What the adjective excludes is passive consumption: the music played in the background of other activity, the diversion that fills time without engaging attention.

\paragraph{Ground projects.} \citet{Williams1981a} argued that a person's ground projects, the commitments that give her reason to go on living, cannot be required of her to be abandoned by an impartial morality without loss of her integrity. The projects that give life its point in Williams's sense are, in the terms of Definition~1, the objects of the deepest engagement, and EXI-03-a uses Williams's related notion of categorical desire, a desire that gives one reason to go on living, and identifies the categorical desires that give life its point with desires for projects one takes to be worthwhile (its Definition~2.1). Engagement, on the conative reading, is the practical form of such desires. Williams's argument against the desirability of immortality \citep{Williams1973}, which EXI-03-a reconstructs, then becomes an argument about the persistence of engagement: an endless life would eventually exhaust the projects with which a person of a given character can engage (Theorem~3.3 of EXI-03-a). EXI-19-b draws the consequence that finitude is a condition of a life's remaining meaningful throughout (its Proposition~1.3).

\lem{2.6}{with projects}

\paragraph{What a project is.} \enquote{Project} has a narrow and a broad sense. In the narrow sense, a project is an extended undertaking with an aim and a completion condition: writing a book, building a house, raising money for a cause. In the broad sense, a project is anything to which a person devotes sustained, directed activity: a relationship, a practice, an inquiry, the care of a garden. Williams's ground projects are projects in the broad sense, since they include personal relationships, and so are Wolf's examples, which include relationships of love and the pursuit of understanding \citep{Wolf2010}.

The text does not say which sense it intends, and the difference matters. On the narrow reading, Definition~1 excludes meaning that comes from love of a person, since a person is not an undertaking, and Assumption~3 becomes doubtful for agents whose circumstances leave no room for undertakings. On the broad reading, Definition~1 includes relationships, and Assumption~3 becomes much easier to defend.

\begin{proposition}[The broad reading]\label{p:broad}
\enquote{Projects} in Definition~1 is to be read in the broad sense: any object of sustained, caring activity is a project in the sense of the definition.
\end{proposition}

\begin{proof}[Argument]
Three considerations require the broad reading. First, the definition is cited to Wolf, whose central cases of meaning include relationships of love, and a definition that excluded them would not be the definition Wolf defends. Second, the narrow reading would make Definition~1 false: lives devoted to the love and care of particular persons are among the paradigms of meaningful lives, and on the narrow reading they would have no meaning in respect of that love. Third, the family of EXI-01 adopts the broad reading. EXI-01-c proves that sustained, caring activity directed at a worthwhile object of engagement is a project of objective worth in the sense of Definition~1 (its Lemma~2.3), and EXI-17-a counts persons, practices and inquiries among the objects of engagement (its Definition~2.2). A reading of a text that makes it false and inconsistent with its family, where another reading makes it true and consistent, is not the correct reading.
\end{proof}

On the broad reading, the later formulations of EXI-17-a and EXI-11-a, which speak of engagement \enquote{with what is worthwhile}, say the same as Definition~1. The word \enquote{projects} contributes something nonetheless: it marks that the engagement must be sustained and directed, not momentary. A single encounter with a worthwhile object, attended to once, is not a project; a relationship of care sustained over time is. EXI-01-c makes the same point in reply to the objection that its lemma trivializes the notion of a project.

\paragraph{The plural.} The definition speaks of projects in the plural and imposes no condition on how they are related. A life may be meaningful through several unrelated engagements, and its meaning is not reduced by their failure to form a whole. This distinguishes Definition~1 from views that require a unifying purpose. \citet{Cottingham2003} argues that a meaningful life requires that one's projects be integrated within a framework that gives them a common point, and that such a framework is in the end religious. EXI-11-a grants that a theistic framework would add meaning, in part by unifying a person's projects under a single purpose that is itself good (its Proposition~4.1), but denies that unification is required. Definition~1 is consistent with both claims: it does not require unity, and its measure is open to the view that a unified set of engagements constitutes a greater extent of engagement than a scattered one.

\paragraph{Projects and defining projects.} \enquote{Project} is used in a different sense in EXI-09, which holds with Sartre that for human beings existence precedes essence. EXI-09-a defines a being's defining project as the end that it adopts and in the light of which it understands itself (its Definition~2.1). A defining project in this sense need not be worthwhile, and a life can have one without being meaningful; EXI-09-b notes that a life can be less meaningful, in the sense of EXI-01, than it might have been, because it engages little with what is worthwhile, without being a defective specimen of anything. The projects of Definition~1 are the objects of a person's engagement; the defining project of EXI-09 is the self-understanding that organizes her ends. The two may coincide, as when a person understands herself through her commitment to a worthwhile cause, but they are different notions.

\lem{2.7}{of objective worth}

\paragraph{Why a qualifier is needed.} The qualifier \enquote{of objective worth} is what makes Definition~1 a hybrid view rather than a subjective one. Without it, the definition would say that a life is meaningful to the extent that it consists in active engagement with projects, whatever their worth, and it would be a version of the view of \citet{Taylor1970}: meaning comes from within, from the will to do what one does. Three cases show that the qualifier is needed.

The first is Taylor's own. Sisyphus is condemned to roll a stone to the top of a hill, from which it rolls down again, forever. Taylor imagines that the gods, in a perverse act of mercy, implant in Sisyphus a compulsive desire to roll stones, so that he does what he most wants to do. Taylor concluded that Sisyphus's life then has all the meaning a life can have. The second is Rawls's grass counter, a person whose only pleasure is to count blades of grass in various geometrically shaped areas such as park squares and well-trimmed lawns \citep[\S65]{Rawls1971}. Rawls introduced him to show that the good of a rational person may consist in odd activities; he has served the literature on meaning as the paradigm of engagement with a worthless object. The third is a case of the kind \citet{Wolf2010} discusses: a person who devotes herself with passion and skill to a pursuit that has no value beyond her attachment to it.

Reflection judges that these lives are not meaningful, or not meaningful in respect of the activities described, although their subjects are actively engaged. The best explanation of the judgment is that engagement confers meaning only when its object has worth that the engagement does not confer. That is what the qualifier says. Whether the explanation is correct, and whether there are any objects with worth of the required kind, is the question of Assumption~2, and the defence in Section~\ref{s:defence} returns to these cases.

\paragraph{\enquote{Worth}.} The text says \enquote{worth}, not \enquote{value}, and Assumption~2 glosses the first in terms of the second (\Lm{3.2}). The words are close, but \enquote{worth} carries an implication that \enquote{value} lacks: that what has it is worthy of something, here of engagement. A project of objective worth is one that merits the engagement it receives. This links Definition~1 to the fitting-attitude tradition, in which values are understood in terms of the attitudes that are fitting toward their bearers, and to EXI-04-a's principle that an attitude is fitting to the extent that the evaluative features under which it presents its object are features the object has (its Principle~P1). Active engagement presents its object as worth engaging with. It is fitting when the object is, and Definition~1 says that only fitting engagement confers meaning. This is why Wolf calls her view the fitting-fulfilment view.

\paragraph{The kinds of worth.} Definition~1 does not say what has objective worth, and the text gives no list. ETH-01 holds that the good has seven irreducible components: well-being, knowledge, achievement, relationships, fairness, beauty and autonomy. EXI-17-a takes these as examples of the worthwhile (its comment on Definition~2.2), and the natural reading of Definition~1 within the collection does the same: projects of objective worth are projects that serve or realize one or more of these goods, the well-being of others, understanding, the making or appreciation of beautiful things, relationships of love and friendship, the advancement of fairness. Two features of this list matter for the definition. It is plural, so that meaning can be found in many different kinds of project, which is what the availability claim needs. And it is not moral: worthwhile projects include aesthetic and epistemic projects that are morally neutral, so that meaning is not the same as moral goodness. A life devoted to mathematics can be meaningful without being morally admirable, and a morally decent life can lack meaning if it engages with nothing.

\paragraph{The degree of worth.} \enquote{Of objective worth} is grammatically a threshold: projects either have objective worth or lack it. But worth comes in degrees, and EXI-11-a states as a principle that engagement with what is of greater worth makes a life more meaningful, other things equal (its Principle~P2). The comment on \Lm{2.3} has argued that the measure in Definition~1 must be sensitive to the worth of the objects of engagement. The threshold reading and the degree reading are therefore both correct, at different points: a project contributes meaning only if it has some objective worth, and it contributes more meaning the greater its worth.

\lem{2.8}{\citep{Wolf2010}}

\paragraph{What is taken from Wolf.} The citation acknowledges that Definition~1 is Wolf's view. Wolf's lectures argue for three claims that the definition presupposes: that meaning is a distinct value, not reducible to happiness or to morality; that it has a subjective and an objective component, neither of which is sufficient; and that the objective component is worth that does not derive from the subject's attitudes. Her slogan, quoted in EXI-20-a, is that meaning arises when subjective attraction meets objective attractiveness. Definition~1 follows closely Wolf's own formula, in her earlier paper, that meaningful lives are lives of active engagement in projects of worth \citep{Wolf1997}.

\paragraph{What is added.} Definition~1 adds three things to Wolf's formulation. It makes the degree explicit, as a variable, where Wolf speaks of lives as meaningful \enquote{insofar as} they are engaged. It replaces Wolf's emphasis on fulfilment, a word that suggests a feeling, by \enquote{active engagement}, which the comment on \Lm{2.5} has read conatively; Wolf's own descriptions of engagement include affective terms, and the conative reading is a choice the commentary has defended, not a report of Wolf's view. And it is combined, in Assumption~2, with a definition of \enquote{objective} that is more restricted than the one Wolf's lectures suggest (\Lm{3.4}).

\paragraph{What is not taken.} Wolf's lectures also argue that meaning gives us reasons for action that are neither self-interested nor moral, so that a picture of practical reason divided between egoism and impartial morality leaves out a large part of what moves us. Definition~1 is neutral on this claim. It says what meaning is, not what reasons it gives. Wolf also declines to give a theory of objective worth, holding that our judgments of worth are fallible and revisable and that no theory of them is needed for the account of meaning to be correct. EXI-01-a follows her in giving no theory, but it is not neutral on the status of worth, since Assumption~2 commits it to the existence of objective worth in a specified sense. The collection supplies the theory Wolf declined to give, in ETH-01 and MTE-01, and Section~\ref{s:defence} examines how much of that theory the text needs.

\paragraph{Questions pressed against Wolf.} Three questions that have been pressed against Wolf's view, in the exchange published with her lectures and in the later literature, bear on Definition~1. They are whether the objective component can be specified without either elitism or emptiness, whether the subjective component asks too much of people whose circumstances do not allow absorption in their work, and whether the account matches what empirical psychology finds about the sources of a sense of meaning. The first question is the elitism objection, which EXI-01-d answers; the second is answered by the conative reading of \Lm{2.5} and by EXI-01-c; the third concerns the sense of meaning, which Definition~1 does not define, rather than meaning itself, and is considered at \Lm{2.4}, where we noted that a subject can be mistaken about the meaning of her life.

\subsection{Definition 1 and its rivals: the cases}\label{s:rivals}

The comments on the phrases of Definition~1 have tested it against cases one at a time. We now bring the cases together and compare the verdicts of Definition~1 with those of its principal rivals. The comparison is a commentary device of a familiar kind: it shows what each phrase of the definition does by showing what goes wrong when it is removed or replaced.

\paragraph{The rival schemata.} We state each rival in a simplified schematic form, sufficient to derive its verdicts on the cases. The schemata are ours; they capture the central condition of each family of views and are not offered as complete statements of any author's theory.

\begin{enumerate}[label=(R\arabic*)]
\item \emph{Subjectivism.} A life is meaningful to the extent that its subject does what she most wants to do, or is engaged with what she cares about, whatever its object. Taylor's conclusion about his Sisyphus is the classical statement \citep{Taylor1970}.
\item \emph{Pure objectivism.} A life is meaningful to the extent that its subject's activity is directed at, or realizes, objective value, whatever her attitude toward it.
\item \emph{The good cause schema.} A life is meaningful to the extent that it causes good outcomes, whether or not they are intended or known; this is the central condition of the account of \citet{Smuts2013}.
\item \emph{The transcendence schema.} A life is meaningful to the extent that it makes a difference to, and leaves traces in, something of value beyond itself; this is one central condition of the view of \citet[ch.~6]{Nozick1981}.
\item \emph{The purpose schema.} A life is meaningful only if it fulfils a purpose assigned to it by a creator; this is the purpose theory of EXI-11-a (its Definition~2.3).
\item \emph{Definition~1.} A life is meaningful to the extent that it consists in active engagement, read conatively, with projects of objective worth, read in the broad sense.
\end{enumerate}

\paragraph{The cases.} The cases are those that the comments have used, together with two from the family. In each, the question is whether the life is meaningful in respect of the activity described.

\begin{enumerate}[label=(c\arabic*)]
\item Taylor's Sisyphus, given a compulsive desire to roll stones (\Lm{2.7}).
\item Rawls's grass counter (\Lm{2.7}).
\item The mathematician who never solves her problem (EXI-17-a, Lemma~3.1).
\item The nurse on the ward where all patients die (EXI-17-a, Lemma~3.1).
\item The listener who studies music and produces nothing (EXI-17-a, Lemma~3.2).
\item The idle inventor, whose joke device later saves many lives (\Lm{2.4}).
\item The depressed nurse, who works with care but without feeling (\Lm{2.5}).
\item Ivan Ilyich, absorbed in his own advancement (Section~\ref{s:defence}).
\item The prisoner who cares for a fellow prisoner (EXI-01-c).
\item The life that fulfils a degrading purpose assigned by its creator (EXI-11-a, Lemma~3.1).
\end{enumerate}

\begin{center}\small
\begin{tabular}{@{}l c c c c c c c c c c@{}}
\toprule
 & c1 & c2 & c3 & c4 & c5 & c6 & c7 & c8 & c9 & c10 \\
\midrule
Reflective judgment & -- & -- & + & + & + & -- & + & -- & + & -- \\
\midrule
(R1) Subjectivism & + & + & + & + & + & -- & -- & + & + & ? \\
(R2) Pure objectivism & -- & -- & + & + & + & + & + & -- & + & -- \\
(R3) Good cause & -- & -- & -- & + & -- & + & + & -- & + & ? \\
(R4) Transcendence & -- & -- & -- & + & -- & + & + & -- & + & ? \\
(R5) Purpose & ? & ? & ? & ? & ? & ? & ? & ? & ? & + \\
(R6) Definition~1 & -- & -- & + & + & + & -- & + & -- & + & -- \\
\bottomrule
\end{tabular}

\smallskip
{\footnotesize + meaningful in respect of the activity; -- not meaningful in that respect; ? the schema gives no verdict without further facts.}
\end{center}

\paragraph{How the verdicts are derived.} The row of reflective judgments records the verdicts that the discipline accepts and that the Commentary has defended at the lemmata cited: EXI-17-a for c3--c5, the comments on \Lm{2.4}, \Lm{2.5} and \Lm{2.7} for c1, c2, c6 and c7, Section~\ref{s:defence} for c8, EXI-01-c for c9 and EXI-11-a for c10. The other rows are derived from the schemata. Under (R1), the depressed nurse is not doing what gives her satisfaction and is not, on the affective reading the schema requires, engaged; the inventor does not care about his device; Ivan Ilyich is doing what he wants. Under (R2), activity directed at or realizing value suffices, so the inventor and the depressed nurse qualify, while Sisyphus, the grass counter and Ivan Ilyich's advancement do not, since their objects have no value. Under (R3), what matters is good caused: the unsuccessful mathematician and the listener cause little or none, while the inventor causes much and the nurse on the second ward causes the relief of her patients' suffering. Under (R4), a difference made beyond oneself is required: the nurse on the second ward makes a difference to how her patients die, the listener and the unsuccessful mathematician leave no traces in the relevant sense. (R5) gives verdicts only given facts about a creator's purposes, except in c10, which is constructed so that the purpose is fulfilled. (R6) gives the verdicts derived in the comments.

\begin{proposition}[Fit with the cases]\label{p:fit}
Of the six schemata, only Definition~1 agrees with reflective judgment on all ten cases.
\end{proposition}

\begin{proof}
By inspection of the table. Subjectivism disagrees on c1, c2, c7 and c8; pure objectivism on c6; the good cause schema on c3, c5 and c6; the transcendence schema on c3, c5 and c6; the purpose schema on c10, and gives no verdict on the rest. Definition~1 agrees on every case.
\end{proof}

The proposition is not a proof that Definition~1 is true. The cases were chosen by the discipline, and some were chosen to support Definition~1; a defender of a rival can contest the reflective verdicts, as Taylor contested the verdict on c1 and as a defender of the good cause account may contest the verdict on c6. What the table shows is the structure of the dispute. Each rival departs from reflective judgment in a characteristic place: subjectivism where engagement meets worthlessness, objectivism where worth meets indifference, the consequence-based schemata where engagement meets failure. Definition~1 is the conjunction of the two conditions whose separate failures the table records, and the cases on which the rivals fail are, one by one, the cases that show why each conjunct is needed.

Two further observations. First, the only case on which pure objectivism fails is c6, the idle inventor, and it is the case on which reflective judgment is least secure; a reader who doubts the verdict there will find pure objectivism nearly as well supported by the table as Definition~1. The comment on \Lm{2.5} gave the reason for the verdict, and the depressed nurse (c7) shows that Definition~1 on the conative reading is not far from objectivism: it requires care, not feeling. Second, the shape-based views of \citet{Kauppinen2012} and of EXI-19 do not appear in the table, because they do not give verdicts on the meaning of a life in respect of an activity; they concern the meaning of a life as a whole. Objection~\ref{o:shape} in Section~\ref{s:whole} considers them.
